\documentclass[10pt,a4paper]{amsart}
\usepackage[margin=19mm]{geometry}
\usepackage{amsmath,amssymb,mathtools}
\usepackage[scr=rsfs,cal=euler]{mathalfa}
\usepackage{xfrac}
\usepackage[unicode,hidelinks,bookmarks=false]{hyperref}
\newcommand{\ie}{i.e.\ }
\newcommand{\cf}{cf.\ }
\newcommand{\resp}{respectively\ }
\newcommand{\Subsection}{\subsection}
\providecommand{\nobreakdash}{\nobreak\hbox{-}}
\setlength{\emergencystretch}{2em}
\multlinegap0pt
\usepackage{bm}
\usepackage{bbm}
\usepackage{dsfont}
\usepackage{xspace}
\usepackage{cancel}
\usepackage{mathtools}
\usepackage{amsthm}
\makeatletter
\@ifundefined{theorem}{\newtheorem{theorem}{Theorem}}{}
\@ifundefined{lemma}{\newtheorem{lemma}[theorem]{Lemma}}{}
\@ifundefined{proposition}{\newtheorem{proposition}[theorem]{Proposition}}{}
\makeatother
\DeclareMathOperator{\ord}{ord}
\newcommand{\Q}{\mathbf{Q}}
\newcommand{\Z}{\mathbf{Z}}
\title[Order drop, Hecke descent, and a mod $p^4$ supercongruence f]{Order drop, Hecke descent, and a mod $p^4$ supercongruence for symmetric-cube hypergeometric coefficients}
\author{Alex Shvets}
\date{Source: arXiv:2604.06238v2 (CC BY 4.0)}
\begin{document}
\maketitle

\noindent\emph{Source and licence.} Order drop, Hecke descent, and a mod $p^4$ supercongruence for symmetric-cube hypergeometric coefficients. Version arXiv:2604.06238v2 (CC BY 4.0), distributed under the Creative Commons Attribution 4.0 International licence, as recorded in the arXiv licence metadata for this exact version and shown on the abstract page https://arxiv.org/abs/2604.06238v2. Full rights notice: licenses/arxiv-2604.06238-CC-BY-4.0.txt.\par\smallskip

\noindent\textit{Editorial scope.} Counted unit: the Theorem whose source label is \texttt{thm:fricke-vanishing} in the article (source0.tex lines 1264--1273, source label \texttt{thm:fricke-vanishing}), delivered together with the article's own complete original proof of it (source0.tex lines 1275--1452, one \texttt{proof} environment of 178 lines). No mathematics is written, repaired or re-derived: statement and proof are the article's own text line for line. Required context retained uncounted: prop:defect-congruence (lines 689--695); prop:defect-pole-order (lines 707--712); prop:fricke-C (lines 743--776); prop:rr-decomposition (lines 1011--1017); lem:fricke-hecke (lines 1036--1045); lem:ordinf-Tp (lines 1215--1225). Every definition, display and label of those lines is retained unchanged. Omitted from this package and not redistributed: 1--688, 696--706, 713--742, 777--1010, 1018--1035, 1046--1214, 1226--1263, 1274--1274, 1453--1518. Nothing inside a retained excerpt is omitted or altered: every retained body is delivered exactly as the article's lines are written, so no omission receipt of any kind accompanies this package. Citations: the retained excerpts carry no citation command at all, so no bibliography entry of the article is transcribed and no reference is redistributed. Reference adaptation: none was needed and none was made; every reference inside the delivered unit resolves inside it, so no unresolved reference or citation is produced. Printed slips kept literal: no printed word, symbol, hypothesis or number of the counted statement or of its proof was altered. Layout and class: the article's own class is \texttt{amsart} with packages including amsmath, amssymb, amsthm, mathtools, hyperref, array, booktabs, graphicx; the wrapper is the lane's pinned \texttt{amsart} editorial wrapper at 10pt on A4, which restates the theorem-like environments the retained text needs and adds the article's own macro definitions that the retained text needs (\texttt{\textbackslash Q}, \texttt{\textbackslash Z}, \texttt{\textbackslash ord}), with extra wrapper packages (bm, bbm, dsfont, xspace, cancel, mathtools). The delivered numbering is the unit's own. Assets: no retained span contains a figure environment, a graphics-inclusion command, a PDF-page inclusion, a table or a TikZ picture; no image file is copied into this package, the package declares no asset dependency and no image file is redistributed with it. Licence: version arXiv:2604.06238v2 of this article is distributed under the Creative Commons Attribution 4.0 International licence, as recorded in the arXiv licence metadata for this exact version and shown on the abstract page https://arxiv.org/abs/2604.06238v2. A full rights notice is delivered at licenses/arxiv-2604.06238-CC-BY-4.0.txt. Characters: every retained excerpt is pure ASCII, so the delivered engine, XeLaTeX with its default Latin Modern text font, reports no missing character.
\begin{proposition}\label{prop:defect-congruence}
For every $r\ge1$ one has
\[
F_r(q)\equiv \widetilde G_r(q)\pmod{p^4}.
\]
Consequently, $F_r\bmod p^4$ is represented by a weakly holomorphic modular form of weight $5$ and character $\chi_3$ on $\Gamma_0(3)$.
\end{proposition}

\begin{proposition}\label{prop:defect-pole-order}
For every $r\ge1$, the defect $F_r(q)$ has pole order at most $r-1$ at the cusp $\infty$:
\[
F_r(q)=O(q^{-r+1}).
\]
\end{proposition}

\begin{proposition}\label{prop:fricke-C}
Let
\[
w_3:=\begin{pmatrix}0&-1\\ 3&0\end{pmatrix},
\qquad
W_3f:=f|_k w_3.
\]
Then
\[
t|_0W_3=\frac{3^{-6}}{t},
\qquad
\Theta|_3W_3=27i\,t\Theta,
\qquad
\left(\frac{Dt}{t}\right)\!\Big|_2W_3=-\frac{Dt}{t}.
\]
Consequently,
\[
C|W_3=-27i\,tC.
\]
Since $W_3^2=-\mathrm{id}$ on weight $5$, one also has
\[
(tC)|W_3=-\frac{i}{27}C.
\]
Therefore
\[
\ord_0(C)=1,
\qquad
\ord_0(tC)=0.
\]
In particular, for every integer $j\ge0$,
\[
\ord_0\!\left(\frac{C}{t^j}\right)=j+1.
\]
\end{proposition}

\begin{proposition}\label{prop:rr-decomposition}
For every $r\ge1$, the exact defect $\widetilde G_r$ admits a finite expansion
\[
\widetilde G_r=\beta_{-1}tC+\sum_{j\ge0}\alpha_j\frac{C}{t^j}
\]
with coefficients $\beta_{-1},\alpha_j\in\Q$ and only finitely many nonzero $\alpha_j$.
\end{proposition}

\begin{lemma}[Fricke--Hecke intertwining]\label{lem:fricke-hecke}
Let $k\ge0$ and let $p\ge5$ be prime. On $M_k^!(\Gamma_0(3),\chi_3)$ one has
\[
T_pW_3=\chi_3(p)\,W_3T_p.
\]
Equivalently, for every $f\in M_k^!(\Gamma_0(3),\chi_3)$,
\[
T_p(f|W_3)=\chi_3(p)\,(T_pf)|W_3.
\]
\end{lemma}

\begin{lemma}\label{lem:ordinf-Tp}
Let $k\ge0$, let $p\ge5$ be prime, and let
\[
g(q)=\sum_{n\ge N}a_n q^n\in M_k^!(\Gamma_0(3),\chi_3),
\qquad N\ge 0.
\]
Then
\[
\ord_\infty(T_pg)\ge \left\lceil\frac{N}{p}\right\rceil.
\]
\end{lemma}

\begin{theorem}[Vanishing of $F_r$ modulo $p^4$ for all $r\ge1$]\label{thm:fricke-vanishing}
For every prime $p\ge5$ and every integer $r\ge1$ one has
\[
F_r(q)=\Lambda_p\!\left(\frac{C(q)}{t(q)^{rp}}\right)-\frac{C(q)}{t(q)^r}\equiv0\pmod{p^4}.
\]
Equivalently,
\[
\widetilde G_r(q)\equiv0\pmod{p^4}.
\]
\end{theorem}

\begin{proof}
By Proposition~\ref{prop:rr-decomposition} there is a finite expansion
\[
\widetilde G_r=\beta_{-1}tC+\sum_{j=0}^{J}\alpha_j\frac{C}{t^j}
\]
for some $J\ge0$.

\emph{Step 1.} This is the required cusp-adapted decomposition.

\emph{Step 2.} By Proposition~\ref{prop:defect-congruence},
\[
F_r\equiv \widetilde G_r \pmod{p^4},
\]
and in fact
\[
\widetilde G_r-F_r=\chi_3(p)p^4V_p\!\left(\frac{C}{t^{rp}}\right),
\]
whose coefficients lie in $p^4\Z_{(p)}$. Thus the $q$-expansion of $\widetilde G_r$ is $p$-integral.

We record the integrality of the basis elements explicitly. Since $C(q)=3E_{5,\chi_0,\chi_3}(q)\in\Z[[q]]$
and
\[
\frac{1}{t(q)}\in q^{-1}\Z[[q]],
\qquad
H(q)=\frac{q}{t(q)}\in 1+q\Z[[q]],
\]
the basis elements
\[
tC=q+O(q^2),\qquad C=1+O(q),\qquad \frac{C}{t^j}=q^{-j}CH^j\in q^{-j}\Z[[q]]
\qquad (j\ge1)
\]
all have $\Z$-coefficients in their $q$-expansions; in particular they lie in $\Z_{(p)}$.
Combined with the $p$-integrality of $\widetilde G_r$, this shows that in the expansion
\[
\widetilde G_r=\beta_{-1}tC+\sum_{j=0}^{J}\alpha_j\frac{C}{t^j}
\]
one has $\beta_{-1},\alpha_j\in\Z_{(p)}$: indeed, the basis is lower-triangular with respect to the
principal part at $\infty$ (with diagonal $1$), so the coefficients are uniquely recovered from the
$q$-expansion of $\widetilde G_r$ by descending triangular solve over $\Z_{(p)}$.

By Proposition~\ref{prop:defect-pole-order},
\[
F_r=O(q^{-r+1}),
\]
so the coefficients of $q^{-J},q^{-J+1},\ldots,q^{-r}$ in $F_r$ all vanish. Combined with
\[
\widetilde G_r-F_r=\chi_3(p)p^4V_p\!\left(\frac{C}{t^{rp}}\right)\in p^4\Z_{(p)}[[q,q^{-1}]],
\]
this gives
\[
[q^{-n}]\widetilde G_r\equiv 0\pmod{p^4}
\qquad (r\le n\le J).
\]

We now prove $\alpha_j\equiv0\pmod{p^4}$ for all $j\ge r$ by descending induction on $j$.

If $J<r$, the claim is vacuous; in that case every $\alpha_j$ with $j\ge r$ is zero by the expansion
itself and there is nothing to prove. We assume henceforth $J\ge r$.

\emph{Base case $j=J$.} The only basis element contributing to $[q^{-J}]\widetilde G_r$ is $C/t^J$
(other elements $C/t^{j'}$ with $j'<J$ start at $q^{-j'}$ with $j'<J$, and $tC$, $C$ start at
$q^0$ or higher). Hence
\[
[q^{-J}]\widetilde G_r=\alpha_J\cdot[q^{-J}](C/t^J)=\alpha_J\cdot 1=\alpha_J.
\]
Since $J\ge r$, the left side is $\equiv0\pmod{p^4}$, so $\alpha_J\equiv0\pmod{p^4}$.

\emph{Inductive step.} Suppose $\alpha_J,\alpha_{J-1},\ldots,\alpha_{j+1}\equiv0\pmod{p^4}$ for some
$j$ with $r\le j<J$. The coefficient of $q^{-j}$ in $\widetilde G_r$ is
\[
[q^{-j}]\widetilde G_r
=
\alpha_j+\sum_{j'=j+1}^{J}\alpha_{j'}\cdot[q^{-j}](C/t^{j'}),
\]
where each $[q^{-j}](C/t^{j'})\in\Z$. By inductive hypothesis, every $\alpha_{j'}$ with $j'>j$ is
$\equiv0\pmod{p^4}$, so the sum is $\equiv0\pmod{p^4}$. Since $[q^{-j}]\widetilde G_r\equiv0\pmod{p^4}$
(because $j\ge r$), we conclude $\alpha_j\equiv0\pmod{p^4}$.

This proves $\alpha_j\equiv0\pmod{p^4}$ for all $j\ge r$.

\emph{Step 3.} Apply $W_3$ to $C/t^{rp}$. By Proposition~\ref{prop:fricke-C},
\[
\left(\frac{C}{t^{rp}}\right)\Big|W_3
=
-27i\,3^{6rp}Ct^{rp+1},
\qquad
\left(\frac{C}{t^r}\right)\Big|W_3
=
-27i\,3^{6r}Ct^{r+1}.
\]
The first form has order $rp+1$ at $\infty$, so Lemma~\ref{lem:ordinf-Tp} gives
\[
\ord_\infty\!\left(
T_p\!\left(\left(\frac{C}{t^{rp}}\right)\Big|W_3\right)
\right)
\ge
\left\lceil\frac{rp+1}{p}\right\rceil
=
r+1.
\]
By Lemma~\ref{lem:fricke-hecke},
\[
\left(T_p\!\left(\frac{C}{t^{rp}}\right)\right)\Big|W_3
=
\chi_3(p)\,
T_p\!\left(\left(\frac{C}{t^{rp}}\right)\Big|W_3\right),
\]
hence
\[
\ord_\infty\!\left(
\left(T_p\!\left(\frac{C}{t^{rp}}\right)\right)\Big|W_3
\right)\ge r+1.
\]
The second transformed term also has order $r+1$ at $\infty$, so
\[
\ord_\infty(\widetilde G_r|W_3)\ge r+1.
\]

Now
\begin{align*}
(Ct^{-j})|W_3
&=
(C|W_3)\,(t^{-j}|_0W_3) \\
&=
(-27itC)\,(3^{6j}t^j) \\
&=
-27i\,3^{6j}Ct^{j+1}
\qquad (j\ge0).
\end{align*}
and the same formula also covers $j=-1$ because
\[
(tC)|W_3=-\frac{i}{27}C=-27i\,3^{-6}C.
\]
Hence
\[
\widetilde G_r|W_3
=
-27i\left(
\beta_{-1}3^{-6}C+\sum_{j=0}^{J}\alpha_j3^{6j}Ct^{j+1}
\right).
\]
The terms
\[
C,\ Ct,\ Ct^2,\ \dots
\]
have distinct orders $0,1,2,\dots$ at $\infty$, so they are linearly independent.
Since $\ord_\infty(\widetilde G_r|W_3)\ge r+1$ is an exact bound (not a congruence),
and since the scalars $3^{-6}$ and $3^{6j}$ are nonzero, the coefficients of
$C, Ct, Ct^2, \ldots, Ct^r$ in the expansion of $\widetilde G_r|W_3$ vanish exactly.
Hence
\[
\beta_{-1}=0,
\qquad
\alpha_j=0
\qquad (0\le j\le r-1).
\]
These are \emph{exact} equalities, not merely congruences modulo $p^4$.

Combining Steps~2 and~3 gives
\[
\beta_{-1}=0,
\qquad
\alpha_j=0
\quad (0\le j\le r-1),
\qquad
\alpha_j\equiv0\pmod{p^4}
\quad (j\ge r).
\]
Therefore
\[
\widetilde G_r\equiv0\pmod{p^4}.
\]
Finally, Proposition~\ref{prop:defect-congruence} gives
\[
F_r\equiv \widetilde G_r\equiv0\pmod{p^4},
\]
as claimed.
\end{proof}

\end{document}
